\verseline{}{“既然天上的神明决意倾覆亚细亚的国势}
\verseline{}{和普里阿摩斯无辜的族人，骄傲的}
\verseline{}{伊利乌姆已经倒下，涅普顿的特洛伊全从地面冒着烟，}
\verseline{}{我们便受众神的征兆驱策，去寻找异乡的流亡之地}
\verseline{5}{和荒无人烟的国土。我们就在安坦德罗斯城下，}
\verseline{}{佛里吉亚的伊达山麓建造舰队，}
\verseline{}{不知命运将把我们带向何方，哪里容许我们停驻，}
\verseline{}{一面聚集人手。初夏才刚刚开始，}
\verseline{}{父亲安喀塞斯便命我们向命运张帆；}
\verseline{10}{我流着泪离开祖国的海岸和港湾}
\verseline{}{离开特洛伊曾经矗立的平原；我这流亡者被带向深海，}
\verseline{}{同伴、儿子、家神和伟大的众神与我同行。}
\verseline{}{远方有片战神的土地，辽阔的平原有人耕种，\footnote{“战神的土地”（terra Mauortia）指色雷斯，传统神话中战神玛尔斯的崇拜中心与居所；吕库尔戈斯（见第 14 行）为曾驱逐酒神并遭严惩的色雷斯暴君。后文（3.49—56）揭示，普里阿摩斯曾秘密将幼子波吕多鲁斯托付给色雷斯国王波吕墨斯托耳与黄金；特洛伊败亡后，国王杀人夺金。欧里庇得斯悲剧《赫卡柏》亦以此事为背景。}}
\verseline{}{色雷斯人在那里犁田，凶猛的吕库尔戈斯曾统治那里；}
\verseline{15}{那是特洛伊古老的宾国，双方的家神也曾结盟，}
\verseline{}{在我们的福运尚存之时。我被带到那里，在弯曲的海岸}
\verseline{}{奠下最初的城墙，却在不利的命运中开始这项事业，}
\verseline{}{又从自己的名字造出‘埃涅阿斯人’这个名号。}
\verseline{}{我正向狄俄涅之女、我的母亲和众神献祭，}
\verseline{20}{祈求他们庇佑初创的事业，又在海岸}
\verseline{}{向天神至高的君王宰杀一头光亮的公牛。}
\verseline{}{近旁恰有一座土丘，顶上长着}
\verseline{}{山茱萸灌木和桃金娘，密密的枝条如矛杆森立。}
\verseline{}{我走近，试着从地上拔起一丛青翠的枝木}
\verseline{25}{好用繁叶覆盖祭坛，}
\verseline{}{却看见一桩骇人、说来不可思议的奇迹。}
\verseline{}{第一株树刚被拔起，断裂的根脱离土壤}
\verseline{}{黑色的血滴便从那里渗出，}
\verseline{}{以污血玷染大地。冰冷的恐惧}
\verseline{30}{震撼我的四肢，血液因惊惧凝结。}
\verseline{}{我又接着去拔第二株柔韧的枝条，}
\verseline{}{想彻底探明那深藏的原因：}
\verseline{}{第二株的树皮里也随即流出黑血。}
\verseline{}{我心中反复思量，敬告荒野的宁芙\footnote{“父神格拉狄乌斯”（Gradivus pater）为战神玛尔斯的古老祭礼尊号，意为“进军之神”或“踏步者”，在此作为色雷斯的保护神受到埃涅阿斯的祈求；“荒野的宁芙”则为主管当地山林的仙女。埃涅阿斯在遭遇折枝渗血异象后向二者祈祷，求其化凶为吉。埃涅阿斯拔取山茱萸与桃金娘枝条时，枝上流血、土下传声；这些植物由刺穿波吕多鲁斯身体的长矛生根而成。随着流血、地下声音和长矛生根依次出现，埃涅阿斯才逐步认出土丘下的死者是波吕多鲁斯。}}
\verseline{35}{和统辖色雷斯田野的父神格拉狄乌斯，}
\verseline{}{求他们依礼使这景象转吉，减轻征兆的凶意。}
\verseline{}{可当我用更大的力气抓住第三根矛杆般的枝条}
\verseline{}{双膝抵着顽抗的沙地奋力拔扯——}
\verseline{}{我该说出，还是沉默？——一声可悲的呻吟}
\verseline{40}{从土丘深处传来，应答的话语抵达我的耳边：}
\verseline{}{‘埃涅阿斯，为何撕裂可怜的我？如今我已经入土，放过我吧；}
\verseline{}{别让你虔敬的双手染上罪污。特洛伊生下的我}
\verseline{}{对你并非异乡人，这鲜血也并非从无知觉的树干流出。}
\verseline{}{唉，逃离这残酷的土地，逃离这贪婪的海岸：}
\verseline{45}{因为我是波吕多鲁斯；铁矛的庄稼在这里覆盖了}
\verseline{}{被刺穿的我，长成尖利的投枪。’}
\verseline{}{这时，我的心才被双重的恐惧压住，}
\verseline{}{我惊住，头发直立而声音哽在喉中。}
\verseline{}{从前，不幸的普里阿摩斯把这个波吕多鲁斯和重金一起，\footnote{第15、61行的 hospitium 指特洛伊与色雷斯的宾主关系，出自埃涅阿斯的叙述。波吕墨斯托耳受托庇护波吕多鲁斯，却因黄金杀害他，破坏了宾主间的信义。}}
\verseline{50}{秘密托付给色雷斯国王抚养，}
\verseline{}{那时他已不再信任达耳达诺斯人的武力，}
\verseline{}{又眼看着都城被围困封锁。}
\verseline{}{可一见特乌克罗斯人的国势崩溃、福运退去，}
\verseline{}{那国王便追随阿伽门农的事业和胜利的军队，}
\verseline{55}{撕断一切神圣的约束，杀死波吕多鲁斯，}
\verseline{}{以暴力夺取黄金。该诅咒的黄金饥渴啊，}
\verseline{}{你有什么不能逼迫凡人的心去做？恐惧离开骨髓以后，\footnote{auri sacra fames 意为“该诅咒的黄金饥渴”。第四卷412行，诗人以几乎相同的句式呼告“残酷的爱神”；两处分别写贪金与爱欲对凡人的逼迫。}}
\verseline{}{我先向族中选出的首领和父亲}
\verseline{}{禀明众神的怪兆，询问他们的判断。}
\verseline{60}{所有人心意相同：离开这片罪恶的土地，}
\verseline{}{舍弃受了玷污的宾国，给舰队以南风。}
\verseline{}{于是我们重新为波吕多鲁斯举行葬礼，大片}
\verseline{}{泥土堆上坟丘；祭奠亡灵的祭坛立在那里，}
\verseline{}{罩着深青的圣带和黑色柏枝，满怀哀意，}
\verseline{65}{伊利乌姆的妇女依照习俗环立，披散头发；}
\verseline{}{我们献上盛着温热乳汁、泡沫涌起的双耳杯，}
\verseline{}{又以祭酒盘献上祭牲鲜血，把亡魂安顿于墓中}
\verseline{}{并放声作最后一次呼唤。\footnote{波吕多鲁斯的葬仪包括堆坟、为亡魂设坛、深青圣带、黑柏枝、妇女散发，以及乳汁与祭牲鲜血等祭奠。}}
\verseline{}{随后，大海初次显得可以信任，风也平息}
\verseline{70}{给出宁静的海面，轻柔低响的南风召我们出海，}
\verseline{}{同伴们推船下水，挤满岸边：}
\verseline{}{我们驶离港湾，土地与城市渐渐退远。}
\verseline{}{海中央有座神圣的岛，最受}
\verseline{}{涅瑞伊得斯的母亲和爱琴海的涅普顿宠爱，}
\verseline{75}{仁厚的持弓神曾将那沿海岸四处漂流的岛，\footnote{指提洛岛（Delos）。相传勒托曾在此岛诞下阿波罗与狄安娜，原本随波漂流的提洛随后被阿波罗系定于米科诺斯与吉阿罗斯之间。下文第 124 行诗人改以古异称“俄耳提癸亚”（Ortygia，源自希腊语“鹌鹑”）称呼此岛；第 694 行同名的俄耳提癸亚则位于西西里叙拉古。}}
\verseline{}{绑在高耸的米科诺斯与吉阿罗斯之间，}
\verseline{}{使它固定下来，容人居住，不再畏惧风浪。}
\verseline{}{我被带到那里；这最安宁的岛在安全港中}
\verseline{}{接纳了疲惫的我们。上岸后，我们敬拜阿波罗之城。}
\verseline{80}{阿尼乌斯既是人民的国王，也是福玻斯的祭司，}
\verseline{}{额头缠着圣带和神圣的月桂，}
\verseline{}{前来相迎，认出了老友安喀塞斯。}
\verseline{}{我们握手缔结宾谊，走进他的屋宇。}
\verseline{}{我在那以古老岩石筑成的神殿前敬拜神明：}
\verseline{85}{‘廷布拉的神啊，赐我们一处自己的家；赐疲惫者城墙；}
\verseline{}{赐我们后裔和一座长存的城市；保全特洛伊第二座}
\verseline{}{帕加马，保全逃过达那俄斯人与无情阿喀琉斯的残民。}
\verseline{}{我们应追随谁？你命我们去往何方？在哪里安家？}
\verseline{}{父神啊，赐下征兆，进入我们的心灵。’}
\verseline{90}{我刚说完，万物仿佛骤然震动，}
\verseline{}{神殿的门槛和月桂，四周整座山峦}
\verseline{}{也为之摇撼；内殿敞开，神谕的铜鼎发出轰鸣。}
\verseline{}{我们俯身伏地，一个声音传入耳中：}
\verseline{}{‘坚忍的达耳达诺斯人，你们祖先的根系最初\footnote{阿波罗命特洛伊人寻找“古老母亲”，神谕故意不直说地点。安喀塞斯先把共同祖源追到克里特，后来家神才纠正为达耳达诺斯来自意大利；航程因对祖源的解释错误而转折。}}
\verseline{95}{从哪片土地生出，那片土地就会以丰饶的怀抱}
\verseline{}{迎接你们归来。去寻找古老的母亲吧：}
\verseline{}{在那里，埃涅阿斯的家族将统治一切海岸，}
\verseline{}{他的子孙，以及从那些子孙生下的后代。’}
\verseline{}{福玻斯这样说。巨大的喜悦混着骚动}
\verseline{100}{骤然兴起，人人都在询问那座城究竟在哪里，}
\verseline{}{福玻斯召流浪者回到什么地方？}
\verseline{}{父亲于是翻检古人留下的传说，}
\verseline{}{说道：‘首领们，听吧，认清你们的希望：}
\verseline{}{大海中央有克里特，伟大朱庇特的岛屿；}
\verseline{105}{那里有伊达山，是我们种族的摇篮。}
\verseline{}{岛上居民拥有百座大城和极其丰饶的王国；}
\verseline{}{倘若我对听来的传说记得不差，我们最初的祖先\footnote{特乌克罗斯与达耳达诺斯是特洛伊祖源传统中的两支。安喀塞斯以克里特祖源、伊达山及山母神祭仪等推断“古老母亲”指克里特；第111行写科律班忒斯的铜钹，第131行另称克里特为库瑞忒斯人的海岸。古代传统常把这两群神灵相联，家神后来则将达耳达诺斯祖系指向意大利。}}
\verseline{}{特乌克罗斯正是从那里航抵罗埃特翁海岸，}
\verseline{}{为自己的王国选定居所。那时伊利乌姆与}
\verseline{110}{帕加马的堡垒还未矗立，人们住在低谷。}
\verseline{}{因此才有栖居库柏勒山的母神、科律班忒斯的铜钹}
\verseline{}{和伊达山的圣林；因此她的祭仪严守沉默，}
\verseline{}{成对的狮子也俯首受轭，为女主人拉车。}
\verseline{}{所以动身吧，让我们追随众神命令所指的方向；}
\verseline{115}{安抚风神，去往克诺索斯的王国。}
\verseline{}{航程并不遥远；只要朱庇特相助，}
\verseline{}{第三日的曙光就会让舰队停在克里特海岸。’}
\verseline{}{说完，他在祭坛宰杀应有的祭牲，}
\verseline{}{一头公牛献给涅普顿，一头献给你，美丽的阿波罗，}
\verseline{120}{一只黑羊献给风暴，一只白羊献给吉利的西风。}
\verseline{}{传言飞来：统帅伊多墨纽斯被逐出父祖的王国，\footnote{伊多墨纽斯是《伊利亚特》中的克里特王；传闻他已离乡，使特洛伊人相信克里特可安全定居。第二卷称他归乡后家神“倒转”，第十一卷狄俄墨得斯又把他的流亡列入希腊归航灾难。}}
\verseline{}{已经离去，克里特海岸也被人抛弃，}
\verseline{}{屋宇不再有敌兵，那些空下的住所正等着我们。}
\verseline{}{我们离开俄耳提癸亚的港口，在海上飞驰，}
\verseline{125}{越过群峰上响着狂欢的纳克索斯和青翠的多努萨，}
\verseline{}{俄利阿罗斯、雪白的帕罗斯，以及散布海上的}
\verseline{}{库克拉得斯群岛，穿过岛屿密布的海峡。}
\verseline{}{水手的呼声竞相响起；}
\verseline{}{同伴彼此鼓舞：‘去克里特，去寻找祖先！’}
\verseline{130}{一阵风从船尾升起，送我们前行}
\verseline{}{终于，我们滑向库瑞忒斯人古老的海岸。}
\verseline{}{我满怀热望，着手修筑那座企盼已久的城市，\footnote{新城也被命名为帕加马，像后文布特洛同的“小特洛伊”。瘟疫与歉收立即证明这次选址错误；特洛伊人只得重新解释神谕，再次寻找真正的祖地。}}
\verseline{}{称它为帕加马，又勉励因这名字而欣喜的族人}
\verseline{}{爱护炉火，以屋宇筑起城堡。}
\verseline{135}{船只已几乎全被拖上干燥的海岸；}
\verseline{}{青年忙于婚配和新开出的田野；}
\verseline{}{我正在制定法度、分配住宅；忽然，一场侵蚀肢体的}
\verseline{}{可悲瘟疫因天空一带气息败坏而降临，}
\verseline{}{祸及树木与庄稼，带来致命的一年。}
\verseline{140}{人们或交出可爱的生命，或拖着病弱的}
\verseline{}{身体；天狼星又把田野灼成不毛之地；}
\verseline{}{草木枯槁，染病的庄稼拒绝供给食物。}
\verseline{}{父亲催我再次渡过已经走过的海路}
\verseline{}{去往俄耳提癸亚的神谕所，向福玻斯求他施恩：}
\verseline{145}{问他会让疲惫的处境得到怎样的终结，命我们从哪里}
\verseline{}{试求救助，又该把航程转向何方。}
\verseline{}{夜已深，大地上的生灵都被睡眠掌握：\footnote{第150—151行说家神在梦中显现，第173行又说“那并非睡梦”，并描述自己清楚认出他们的面貌。家神将阿波罗神谕指向意大利，并以达耳达诺斯的祖源说明这片土地与特洛伊人的关系。}}
\verseline{}{我从特洛伊、从城中烈火里带出的}
\verseline{}{众神圣像和佛里吉亚的家神}
\verseline{150}{在梦中仿佛立于我卧榻前，显在眼中}
\verseline{}{清晰地沐浴着满室光辉；一轮满月}
\verseline{}{正从嵌入墙壁的窗中洒下光芒；}
\verseline{}{他们于是这样开口，用这些话减轻我的忧虑：}
\verseline{}{‘阿波罗等你被送到俄耳提癸亚才会说的话，}
\verseline{155}{此刻就在这里宣告。看，他主动派我们来到你的门前。}
\verseline{}{达耳达诺斯的城焚毁以后，我们追随你，也追随你的兵戈，}
\verseline{}{在你的统率下，我们随舰队渡过汹涌的大海，}
\verseline{}{也将把未来的子孙高举到星辰之间，}
\verseline{}{把统治权赐给他们的城市。你要为伟大的后代}
\verseline{160}{准备宏伟的城墙，不要在漫长流亡的劳苦前退却。}
\verseline{}{居所必须更换：提洛的神不曾劝你来这片海岸，}
\verseline{}{阿波罗也不曾命你在克里特定居。}
\verseline{}{有一个地方，希腊人称之为赫斯珀里亚，}
\verseline{}{那是一片古老的土地，武力强盛而泥土肥沃；}
\verseline{165}{俄诺特里亚人曾居住于此；如今有传言说，他们的后裔}
\verseline{}{用他们领袖的名字，将这个民族称为意大利：}
\verseline{}{那里才是我们自己的居所；达耳达诺斯出自那里，}
\verseline{}{父祖伊阿西翁也出自那里，我们的种族从这位始祖发源。\footnote{这是家神向埃涅阿斯说明的祖源：达耳达诺斯与伊阿西翁来自意大利，第170行又命他寻找科律图斯和奥索尼亚的土地。安喀塞斯听到转述后，在第180—181行才承认自己混淆了两支祖系。}}
\verseline{}{起来，把这些不容怀疑的话欣然转告你年迈的父亲，}
\verseline{170}{让他寻找科律图斯和奥索尼亚的土地；}
\verseline{}{朱庇特不许你占据狄克忒的田野。’}
\verseline{}{这样的幻象和神明的声音使我震惊——}
\verseline{}{那并非睡梦，我仿佛当面认出了他们的面貌}
\verseline{}{覆着头巾的发髻和亲临眼前的容颜；}
\verseline{175}{冰冷的汗水顿时流遍全身，}
\verseline{}{我猛然起身，仰掌朝向天空，}
\verseline{}{一面高声祷告，一面在炉火上浇下}
\verseline{}{未受玷污的供物。祭礼完成，我满怀喜悦}
\verseline{}{把一切告知安喀塞斯，依次说明事情的经过。}
\verseline{180}{他认出了我们混合的血脉和两支祖系，\footnote{prolem ambiguam geminosque parentes 指特洛伊人混合的血统与两支祖系：来自克里特的特乌克罗斯，以及来自意大利的达耳达诺斯。第181行以“新的错误”与“古老的地名”相对照。}}
\verseline{}{明白自己又一次误认了古老的地名。}
\verseline{}{他随即回忆：‘儿子，经受伊利乌姆命运磨炼的人啊，}
\verseline{}{只有卡珊德拉曾向我预言这样的遭际。}
\verseline{}{如今我想起，她曾说这些事注定属于我们的种族，}
\verseline{185}{又一次次呼唤赫斯珀里亚，呼唤意大利的王国。}
\verseline{}{但谁会相信特乌克罗斯人终将来到赫斯珀里亚的}
\verseline{}{海岸？那时又有谁肯被女先知卡珊德拉打动？}
\verseline{}{让我们服从福玻斯，听从警示，追随更好的道路。’}
\verseline{}{他说完，我们全都欢欣地服从。}
\verseline{190}{我们又舍弃这个居所，只留下寥寥数人}
\verseline{}{张帆出海，驾着中空的木船越过辽阔水面。}
\verseline{}{船队驶入深海以后，陆地已不再}
\verseline{}{显现，四面都是天空，四面都是海，}
\verseline{}{一团青黑的雨云忽然立在我头顶，}
\verseline{195}{带来黑夜和风暴，波浪在幽暗中耸起寒栗。}
\verseline{}{狂风顷刻卷动海水，巨大的}
\verseline{}{波涛涌起；我们四散，在浩瀚的漩涡中颠簸；}
\verseline{}{雨云裹住白昼，湿暗的夜夺走}
\verseline{}{天空；裂开的云层中闪电接连燃起。}
\verseline{200}{我们被赶离航路，在昏暗的浪涛间漂流。}
\verseline{}{就连帕利努鲁斯也说，天色使他不能分辨昼夜，\footnote{帕利努鲁斯是舰队舵手。此处黑雾使他不能辨认昼夜与航线；第五卷末，他又在睡神的作用下坠海。}}
\verseline{}{身在大海中央，他已经记不得航线。}
\verseline{}{我们就在不明方向的黑雾里迷航三个白日，}
\verseline{}{又经过同样多不见星光的夜晚。}
\verseline{205}{直到第四日，陆地才终于仿佛从海上升起，}
\verseline{}{远远显出山峦，翻卷着烟雾。}
\verseline{}{风帆落下，我们俯身奋桨；水手毫不迟延}
\verseline{}{一齐划开泡沫，扫过青蓝的海面。}
\verseline{}{我从风浪中获救，最先接纳我的}
\verseline{210}{是斯特洛法得斯的海岸；这些岛取了希腊名字，}
\verseline{}{立在辽阔的伊奥尼亚海中，可怕的刻莱诺}
\verseline{}{和其他哈耳庇厄居住在那里。自从菲纽斯的家门}
\verseline{}{被关闭之后，她们便因恐惧舍弃先前的筵席。}
\verseline{}{再没有什么怪物比她们更阴惨，也没有更残酷的}
\verseline{215}{灾害与神怒从斯提克斯的水中升起。\footnote{斯特洛法得斯群岛的哈耳庇厄承接阿波罗尼俄斯《阿尔戈英雄纪》第二卷的菲纽斯故事：哈耳庇厄不断抢夺并污秽菲纽斯的食物，随后遭北风之子追逐。维吉尔把这组怪物移入特洛伊人的航程，并让刻莱诺开口预言。}}
\verseline{}{这些怪鸟长着少女的面孔，腹中}
\verseline{}{泻出最污秽的东西，双手如钩，脸总因}
\verseline{}{饥饿而苍白。}
\verseline{}{我们被带到这里，刚刚驶入港湾，看啊}
\verseline{220}{就见丰美的牛群散在田野各处，}
\verseline{}{还有无人看管的山羊群在草地上游荡。}
\verseline{}{我们持刀扑去，呼唤众神，也请朱庇特本人}
\verseline{}{来分享战利品；随后在弯曲的海岸}
\verseline{}{铺设坐榻，享用丰盛的宴席。}
\verseline{225}{可哈耳庇厄骤然从山上俯冲而下，}
\verseline{}{骇人的翅声大作，}
\verseline{}{她们抢夺食物，以不洁的触碰玷污一切，}
\verseline{}{刺耳的尖叫又从恶臭中传来。}
\verseline{}{我们再次退到远处，躲在一块凹入的岩石下，}
\verseline{230}{周围树林封闭，阴影森然，\footnote{Conte 2019 将第230行以方括号排印；该行重现第一卷311行。该版校勘记注明塞尔维乌斯未释此行，并记录 Ribbeck 的删疑。}}
\verseline{}{重新摆开餐桌，在祭坛恢复火焰：}
\verseline{}{那群喧响的怪鸟又从天空另一边}
\verseline{}{从黑暗的藏处飞出，以钩爪绕着猎物盘旋，}
\verseline{}{用嘴污染食物。我于是命同伴们}
\verseline{235}{拿起兵器，向这可怕的族类开战。}
\verseline{}{他们依令而行，把剑隐在草中，}
\verseline{}{又将盾牌藏好，不露痕迹。}
\verseline{}{她们一俯冲，在弯曲的海岸上响起翅声，}
\verseline{}{米塞努斯就从高处的瞭望台}
\verseline{240}{吹响中空的铜号。同伴们冲上前去，尝试一场新奇的战斗，}
\verseline{}{要以铁器伤害这些海上污秽的怪鸟：}
\verseline{}{可她们的羽毛不受刀剑之力，背上也不接纳}
\verseline{}{任何伤口，转眼飞上星空，}
\verseline{}{只留下吃剩的猎物和污秽的痕迹。}
\verseline{245}{唯有刻莱诺停在一块极高的岩石上，}
\verseline{}{这带来不幸的女先知从胸中迸出这样的声音：}
\verseline{}{‘杀了我们的牛，放倒了小公牛以后，}
\verseline{}{拉俄墨冬的子孙，你们还准备发动战争，}
\verseline{}{把无辜的哈耳庇厄逐出祖传的王国？}
\verseline{250}{那么用心听取，把我的这些话钉牢在心中，}
\verseline{}{全能的父神曾向福玻斯预言，福玻斯·阿波罗又向我}
\verseline{}{预言；我，复仇女神中最强大者，现在为你们揭晓。}
\verseline{}{你们正循航路、借着召来的风驶向意大利；\footnote{Conte 2019 在第253行 uentisque uocatis 后断句，使“航路”与“召来的风”一同连接前面的 Italiam cursu petitis。该版校勘记另记一种句读，将 uentisque uocatis 接到下一句 ibitis Italiam。}}
\verseline{}{你们将抵达意大利，也将获准进入港湾；}
\verseline{255}{但在你们用城墙围住命定赐下的城市以前，}
\verseline{}{可怕的饥饿和你们对我们施暴所结下的罪债}
\verseline{}{将迫使你们用牙齿咬碎桌案，全部吞下。’\footnote{刻莱诺预言饥饿将迫使众人“吃掉桌子”。第七卷他们以面饼托食，吃完食物又吃面饼，尤卢斯玩笑说“连桌子也吃了”；可怕预言届时以日常动作兑现。}}
\verseline{}{她说完，振翼逃回森林。}
\verseline{}{突来的恐惧使同伴们的血液冰冷}
\verseline{260}{凝固；众人心气尽丧，不再想凭兵器求胜，}
\verseline{}{而要求改用誓愿和祈祷来乞求和平，}
\verseline{}{无论她们是女神，还是阴惨污秽的怪鸟。}
\verseline{}{父亲安喀塞斯便站在海岸，伸开双掌，}
\verseline{}{呼求伟大的神明，宣布应有的祭献：}
\verseline{265}{‘众神啊，挡开这威胁；众神啊，移走这样的厄运，}
\verseline{}{以慈和守护敬神之人！’随后他命人从岸边}
\verseline{}{扯下缆绳，抖开松脱的帆索。}
\verseline{}{南风张满船帆；我们从泛沫的波浪上逃走，}
\verseline{}{航向由风和舵手一同指引。}
\verseline{270}{树木葱郁的扎金索斯已在海中显现，}
\verseline{}{还有杜利基乌姆、萨墨和岩石高耸的涅里托斯。}
\verseline{}{我们避开伊塔卡的礁石，避开拉埃耳忒斯的王国，}
\verseline{}{诅咒那片哺育残酷乌利塞斯的土地。\footnote{伊塔卡、拉埃耳忒斯与乌利塞斯把这段航程接入《奥德赛》的地理。特洛伊人避开旧敌的岛屿，继续寻找意大利这一祖源之地，并在那里建立新家园。}}
\verseline{}{不久，琉卡忒斯山云雾笼罩的峰顶}
\verseline{275}{和令水手畏惧的阿波罗神庙也展现出来。}
\verseline{}{疲惫的我们向那里驶去，靠近那座小城；}
\verseline{}{锚从船首抛下，船尾沿岸停列。}
\verseline{}{终于获得这意外来到的陆地，}
\verseline{}{我们向朱庇特行净身礼，在祭坛点火还愿，}
\verseline{280}{又在阿克提乌姆海岸以伊利乌姆的竞技欢庆。}
\verseline{}{同伴们裸身涂上流动的橄榄油，}
\verseline{}{操练祖国的角力；能逃过那么多阿尔戈斯人的}
\verseline{}{城市，从敌阵中央坚持逃亡，令人欣慰。}
\verseline{}{其间太阳绕完一个漫长的年轮，}
\verseline{285}{寒冬以北风吹皱海面。}
\verseline{}{我把一面中空的青铜盾——伟大的阿巴斯曾经佩带——}
\verseline{}{钉在神庙正面的门柱，又以铭文标明此事：}
\verseline{}{‘埃涅阿斯从胜利的达那俄斯人手中取得这些武器。’\footnote{埃涅阿斯在阿克提乌姆举行“伊利乌姆竞技”，又把阿巴斯之盾题作从胜利者达那俄斯人手中所得。战败的特洛伊人此刻在希腊土地上暂时扮演胜利者；阿克提乌姆这一地点也会使奥古斯都时代的读者想到后来决定罗马世界归属的海战。}}
\verseline{}{随后我命大家离开港湾，坐上桨手的长凳：}
\verseline{290}{同伴竞相击水，以船桨扫过海面。}
\verseline{}{我们随即把法埃阿克斯人高入云天的城堡抛在身后，}
\verseline{}{沿伊庇鲁斯海岸航行，进入}
\verseline{}{卡奥尼亚的港湾，来到布特洛同高城。}
\verseline{}{一桩难以置信的传闻在这里传入耳中，}
\verseline{295}{普里阿摩斯之子赫勒努斯正统治希腊人的城市，}
\verseline{}{继承埃阿科斯后裔皮洛士的妻子与权杖，}
\verseline{}{而安德罗玛刻又回到一位同胞丈夫身边。\footnote{赫勒努斯是普里阿摩斯之子和先知；安德罗玛刻是赫克托耳遗孀。她先为皮洛士／涅俄普托勒摩斯所俘，后转交赫勒努斯；俄瑞斯忒斯杀死皮洛士后，赫勒努斯得到部分王国。}}
\verseline{}{我惊住，胸中燃起奇异的渴望，}
\verseline{}{想要和这位故人交谈，了解如此巨大的变故。}
\verseline{300}{我离开舰队和海岸，从港口向前走去，}
\verseline{}{恰逢安德罗玛刻正在奉上周年祭宴与悲哀的供物，}
\verseline{}{地点在城外林中、那条仿造的西摩伊斯河畔；}
\verseline{}{她向亡者的骨灰奠献，呼唤亡魂，}
\verseline{}{来到赫克托耳的坟前——那座空冢覆着青绿草皮}
\verseline{305}{并有两座引泪的祭坛，都是她亲手奉献。\footnote{安德罗玛刻在仿造的西摩伊斯河旁，为赫克托耳设空冢与双坛。第349—351行又写“小特洛伊”、帕加马、仿称克珊托斯的枯溪，以及斯开亚门；这些沿用故城名称的景物共同保存特洛伊记忆。}}
\verseline{}{她看见我走近，又望见周围特洛伊的}
\verseline{}{武装，神志迷乱，被这巨大异象惊吓}
\verseline{}{当场僵住，温热离开骨髓；}
\verseline{}{她倒了下去，过了很久才终于开口：}
\verseline{310}{‘女神之子，你这面容可是真的？你真身来到我面前了吗，}
\verseline{}{你还活着？若滋养生命的光已经离去，}
\verseline{}{那么赫克托耳在哪里？’说罢，她泪水涌出}
\verseline{}{哭声充满整个地方。她近乎疯狂，我只能勉强}
\verseline{}{接上几句话，心神不定，断续出声：}
\verseline{315}{‘我确实活着，拖着生命走过一切绝境；}
\verseline{}{不要怀疑，你所见的是真实之身。}
\verseline{}{唉，你从如此伟大的丈夫身边坠落后，什么遭际}
\verseline{}{接住了你？又有什么足以与你相称的命运再次眷顾你，}
\verseline{}{赫克托耳的安德罗玛刻啊？你仍是皮洛士的妻子？’}
\verseline{320}{她低下眼睛，以压低的声音回答：}
\verseline{}{‘唯有那位普里阿摩斯之女比众人幸运，}
\verseline{}{她奉命死在敌人的坟前、特洛伊高耸的城墙下，}
\verseline{}{从未承受过把她分给任何主人的抽签，}
\verseline{}{也不曾以俘虏之身触碰胜利者的床榻！\footnote{“普里阿摩斯之女”是波吕克塞娜，在阿喀琉斯墓前被杀的故事见欧里庇得斯《赫卡柏》。安德罗玛刻认为她未受抽签分配与主人床榻之辱，因而比其他俘虏幸运。《安德罗玛刻》写她与涅俄普托勒摩斯所生的莫洛索斯；第327行“在奴役中生育”可联系这一背景。}}
\verseline{325}{至于我们，祖国焚毁以后被载过各处海洋，}
\verseline{}{在奴役中生育，忍受阿喀琉斯一脉的傲慢，}
\verseline{}{也忍受那骄横的青年。后来他追求}
\verseline{}{勒达的孙女赫尔弥俄涅和斯巴达的婚姻，}
\verseline{}{便把我这个女奴交给同为奴仆的赫勒努斯占有。}
\verseline{330}{可是俄瑞斯忒斯对被夺去的未婚妻燃烧着强烈爱欲，}
\verseline{}{又受自身罪行的复仇女神催逼，}
\verseline{}{趁皮洛士不备，在他父亲的祭坛旁把他杀死。\footnote{俄瑞斯忒斯原与赫尔弥俄涅订婚，皮洛士却把她娶走；俄瑞斯忒斯又因弑母罪受复仇女神追逐。塞尔维乌斯（3.332）把“父亲的祭坛”一解为祭祀阿喀琉斯的祭坛，也列阿波罗祭坛等异解。按前一解，此死法可与第二卷皮洛士在祭坛前杀害普里阿摩斯相照应。}}
\verseline{}{涅俄普托勒摩斯死后，一部分应得的王国}
\verseline{}{转归赫勒努斯。他取特洛伊人卡翁的名字}
\verseline{335}{把平原称为卡奥尼亚，也让整个国土沿用此名，}
\verseline{}{又在山脊建起一座帕加马，这座伊利乌姆的城堡。}
\verseline{}{但什么风、什么命运把你送上这条航路？}
\verseline{}{又是哪位神把毫不知情的你送到我们海岸？}
\verseline{}{孩子阿斯卡尼乌斯如何？他是否还活着，还呼吸空气？}
\verseline{340}{那个在特洛伊就已为你……的孩子？}
\verseline{}{他是否还记挂失去的母亲？}
\verseline{}{他的父亲埃涅阿斯和舅父赫克托耳，}
\verseline{}{是否激起他祖传的勇武和男子气概？’\footnote{Conte 2019 第340行只存 quem tibi iam Troia…?；该版校勘记引塞尔维乌斯，说明此半行语意未完。第341行“失去的母亲”通常联系克瑞乌萨，塞尔维乌斯（3.341）则以 parentis 指祖国。安德罗玛刻还问阿斯卡尼乌斯是否继承父亲与舅父的勇武；临别赠衣时，她说他的眼睛、手与面容像阿斯提阿那克斯，两人本应同龄。}}
\verseline{}{她流着泪倾泻这些话，徒然唤起}
\verseline{345}{漫长的悲泣；这时英雄赫勒努斯}
\verseline{}{普里阿摩斯的儿子，在众多随从簇拥下从城里出来，}
\verseline{}{认出亲族，满怀喜悦地把我们领向门内，}
\verseline{}{话语之间，他的泪水不断流下。}
\verseline{}{我向前走去，认出一座小小的特洛伊，认出}
\verseline{350}{仿造伟大帕加马的堡垒和一条冠以克珊托斯之名的枯溪，}
\verseline{}{又拥抱斯开亚门的门槛。\footnote{埃涅阿斯认出“小特洛伊”、帕加马和仿称克珊托斯的枯溪，并拥抱斯开亚门。布特洛同沿用故城的名称与地貌，与他尚待建立的新城形成对照。}}
\verseline{}{特乌克罗斯人也和我一起享用这座友城：}
\verseline{}{国王在宽阔的柱廊接待他们；}
\verseline{}{他们在大厅中央奠下巴克科斯的酒，}
\verseline{355}{金器上摆着宴席，人人手执酒碗。}
\verseline{}{两日已经过去，风开始召唤船帆}
\verseline{}{帆布被劲吹的南风鼓起。}
\verseline{}{我以这些话向先知开口，向他询问：}
\verseline{}{‘特洛伊所生、众神的传译者啊，你感知福玻斯的神意，}
\verseline{360}{感知三足鼎、克拉罗斯的月桂和星辰，\footnote{Conte 2019 正文采用 S. Ottaviano 据斯塔提乌斯拟复的 tripodas laurusque Clari；主要传本在 Clari/Clarii 与 laurus 的结构上存在紊乱。此读法把三足鼎与克拉罗斯的月桂并列为阿波罗神谕的标志。}}
\verseline{}{又懂得飞鸟的语言与羽翼疾飞的预兆，}
\verseline{}{来吧，告诉我——种种祭兆都已许诺我的航程}
\verseline{}{将会顺遂，所有神明也以神意劝我}
\verseline{}{前往意大利，探寻那遥远的土地：}
\verseline{365}{唯有哈耳庇厄刻莱诺预言一桩前所未有、说来不祥的}
\verseline{}{怪兆，宣告惨痛的神怒}
\verseline{}{和污秽的饥饿。我首先要躲避哪些危险？}
\verseline{}{又要沿着什么道路才能战胜这重重苦难？’}
\verseline{}{赫勒努斯首先依照习俗宰杀小公牛，}
\verseline{370}{祈求众神息怒，解开神圣额头上的祭带，}
\verseline{}{随后亲手把我领到你的门槛，福玻斯啊，}
\verseline{}{我因强烈的神临而敬畏悬心；祭司接着}
\verseline{}{从神灵感发的口中唱出这些预言：}
\verseline{}{‘女神之子——确凿的证据表明，你在更崇高的}
\verseline{375}{征兆下航过深海。众神之王这样分配命运}
\verseline{}{转动种种变迁，秩序就这样循环——}
\verseline{}{我只从许多事情里解释少数几件，好让你更安全地经过}
\verseline{}{陌生海域，停驻奥索尼亚的港湾，}
\verseline{}{因为帕耳卡三女神不许赫勒努斯知道其余事情，}
\verseline{380}{萨图尔努斯之女朱诺也不准他说出。\footnote{第379—380行，帕耳卡三女神不许赫勒努斯知道其余事情，朱诺则不许他说出：scire 限制知晓，fari 限制言说。}}
\verseline{}{首先，你自以为意大利已近在眼前，}
\verseline{}{无知地准备进入邻近港湾；其实一条漫长的无路之途}
\verseline{}{隔着遥远陆地，把它远远挡在前方。}
\verseline{}{你还须先在特里纳克里亚海水中弯桨，}
\verseline{385}{让船队巡过奥索尼亚的盐海，经过}
\verseline{}{冥界之湖与喀耳刻的埃埃亚岛，此后，}
\verseline{}{才能在安全的土地上建起城市：}
\verseline{}{我会告诉你征兆，你要把它们深藏心中：}
\verseline{}{当你忧心地走到一条隐秘河流的水边，}
\verseline{390}{看见一头巨大的母猪在岸边橡树下}
\verseline{}{产下一窝三十只幼崽，卧在那里。}
\verseline{}{白色的母猪伏在地上，白色的孩子围着乳房，}
\verseline{}{那里便是建城之地，也是劳苦必将止息之处。\footnote{白母猪与三十只白幼崽是建城征兆，第八卷43—45行再次预告，随后在81—85行出现。第八卷47—48行又说，三十年后阿斯卡尼乌斯将建立阿尔巴；三十头幼崽与三十年相呼应，Alba 也与 alba（“白色”）相应。}}
\verseline{}{别因预言说你们将咬食桌案而惊恐：}
\verseline{395}{命运会找到道路，受召的阿波罗也会前来。\footnote{赫勒努斯安慰埃涅阿斯，不必害怕“吃桌子”的预言：命运会找到道路。预言在第七卷以吃尽承托食物的面饼应验。}}
\verseline{}{至于这些土地，以及意大利这段最近的海岸，}
\verseline{}{就在邻接我们海域的浪潮中受海水冲刷，}
\verseline{}{你要避开；那里所有城垒都住着凶恶的希腊人。}
\verseline{}{纳律基翁的洛克里斯人已经在这里修起城墙，}
\verseline{400}{吕克托斯的伊多墨纽斯率军占据萨伦提努斯的平原；}
\verseline{}{墨利波亚统帅菲罗克忒忒斯那座小小的}
\verseline{}{佩特利亚也在那里，凭借城墙自固。}
\verseline{}{而且，等舰队渡海以后停驻，}
\verseline{}{你在岸边设起祭坛，准备履行誓愿，}
\verseline{405}{要以紫色外袍盖住头发，}
\verseline{}{免得在敬奉诸神的圣火之间，}
\verseline{}{出现敌人的面孔，扰乱征兆。}
\verseline{}{你的同伴要守住这项祭仪，你自己也要守住；}
\verseline{}{愿后代在这一敬神传统中永葆纯洁。\footnote{祭祀者以紫袍遮头，使敌人的面孔无法扰乱征兆；这一做法后来被叙述为罗马祭仪传统。诗中把这一习俗追溯到赫勒努斯给埃涅阿斯的航海指令。}}
\verseline{410}{而当风把离开那里以后的你送近西西里海岸，}
\verseline{}{狭窄的佩洛鲁斯屏障渐渐向两边分开，}
\verseline{}{你要在漫长的迂回中寻求左边的陆地}
\verseline{}{和左边的海面，逃开右方的海岸与波涛。}
\verseline{}{据说这些地方从前被猛烈的力量和浩大的崩坏}
\verseline{415}{撕裂——漫长久远的岁月竟能改变如此之多——}
\verseline{}{原先两片土地曾经连成一体}
\verseline{}{大海却以暴力涌进中间，用波涛}
\verseline{}{从西西里切断赫斯珀里亚的一侧}
\verseline{}{狭窄的潮水流在分处两岸的田野与城市之间。}
\verseline{420}{斯库拉盘踞右方，不能平息的卡律布狄斯占据左方，}
\verseline{}{她一天三次用深渊底部的漩涡吞下}
\verseline{}{浩荡的巨浪，落入陡绝的深处，又轮番把海水}
\verseline{}{举向天空，以浪涛鞭打星辰。}
\verseline{}{斯库拉却被洞窟关在黑暗的藏处，}
\verseline{425}{向外伸出一张张巨口，把船只拖向礁石。}
\verseline{}{她上身是人的面貌，是一位胸膛美丽的少女，}
\verseline{}{直到腰腹；下身却是庞然海怪，}
\verseline{}{海豚的尾部接在恶狼的肚腹。}
\verseline{}{最好迟缓地绕过特里纳克里亚的帕基努斯岬，}
\verseline{430}{将航线弯成长远的迂回，}
\verseline{}{也胜过在浩大洞穴之下看见一次}
\verseline{}{丑陋的斯库拉，听那礁石回响青黑恶犬的吠声。\footnote{斯库拉与卡律布狄斯来自《奥德赛》的海峡怪物传统。赫勒努斯要求宁可绕行也不要从两者之间直穿；诗中以两怪分踞左右的身体描出这片危险海峡。不同古代传统对这片海域各有安排。}}
\verseline{}{此外，若赫勒努斯还有一点先见，若这位预言者}
\verseline{}{尚值得信任，若阿波罗以真理充满他的心，}
\verseline{435}{女神之子啊，我只把这一件事置于一切之上}
\verseline{}{预告给你，并要一再、一再地劝诫：}
\verseline{}{首先以祈祷敬拜伟大朱诺的神意；}
\verseline{}{欣然向朱诺许愿，以谦卑的供物}
\verseline{}{赢得强大的女主人；只有这样，你才能最终战胜阻碍}
\verseline{440}{离开特里纳克里亚，被送往意大利的疆界。}
\verseline{}{等你被带到那里，抵达库迈城，}
\verseline{}{来到神圣的湖泊与树林间轰鸣的阿韦尔努斯，}
\verseline{}{你会看到那发狂的女先知；她在岩窟深处}
\verseline{}{歌唱命运，把符号与字句托付给叶片。}
\verseline{445}{少女把写在叶上的每一篇神谕}
\verseline{}{依照次序排好，留在封闭的洞穴中。}
\verseline{}{它们安静地停在原位，不会离开秩序；}
\verseline{}{可一旦门轴转动，门页送入}
\verseline{}{一阵微风，搅乱柔弱的叶片，}
\verseline{450}{她便再不费心去岩洞里捉住那些飞舞的叶子，}
\verseline{}{不把它们放回原处，也不重新拼合神谕：}
\verseline{}{求问者只得未经指点便离开，憎恨西比尔的居所。\footnote{库迈西比尔把神谕写在叶片上；门风吹乱以后，她不再整理。赫勒努斯因此要求埃涅阿斯坚持让她亲口吟唱。第六卷实际进入洞窟时，口头神谕与神灵附体成为中心。}}
\verseline{}{在那里，任何迟延都不要使你觉得损失太大}
\verseline{}{纵然同伴催促，纵然航程强烈地召帆入海}
\verseline{455}{顺风已能把帆腹吹满}
\verseline{}{你仍要走近女先知，以祈祷索求神谕}
\verseline{}{请她亲自歌唱，甘愿开口发声。}
\verseline{}{她会为你解说意大利的部族和将临的战争，}
\verseline{}{说明你该以什么方式躲避或承受每一场苦难，}
\verseline{460}{她一旦受你敬奉，便会赐下顺利的航路。}
\verseline{}{这些就是我的声音获准告诫你的事。}
\verseline{}{去吧，以功业把伟大的特洛伊高举到苍穹。’}
\verseline{}{预言者以友善的话说完这些以后，}
\verseline{}{便命人把嵌满黄金与雕花象牙的沉重礼物}
\verseline{465}{运上船，又往船舱装进}
\verseline{}{大量白银与多多纳的铜鼎，}
\verseline{}{一件钩环相接、以黄金织成三层的锁子甲，}
\verseline{}{一顶以高耸锥顶和浓密盔缨为饰的显赫头盔，}
\verseline{}{涅俄普托勒摩斯的武器；也有赫勒努斯送给我父亲的礼物。}
\verseline{470}{他又添上马匹，添上向导；}
\verseline{}{补足桨手，同时给同伴装备兵器。}
\verseline{}{其间，安喀塞斯正命人给舰队装好风帆}
\verseline{}{免得耽误了送行的风。}
\verseline{}{福玻斯的传译者带着无比敬意对他说：}
\verseline{475}{‘安喀塞斯啊，维纳斯以高贵的婚姻眷顾你，}
\verseline{}{你是众神关怀之人，两次从帕加马的废墟获救，}
\verseline{}{看，奥索尼亚的土地就在眼前；张帆占有它吧。}
\verseline{}{但你仍必须在海上驶过这一段海岸；}
\verseline{}{阿波罗为你开启的奥索尼亚还在遥远的另一边。}
\verseline{480}{走吧，’他说，‘因儿子的敬义而幸福的人啊。我何必\footnote{赫勒努斯称安喀塞斯“因儿子的敬义而幸福”。这里的pietas是埃涅阿斯对父亲、家神和使命的尽责，不只是私人孝顺；这句称赞着眼于他护送父亲与家神的行为。}}
\verseline{}{再往下说，以话语耽误正在升起的南风？’}
\verseline{}{安德罗玛刻也为最后的离别满怀哀伤}
\verseline{}{拿来以金线织出图纹的衣裳}
\verseline{}{和一件给阿斯卡尼乌斯的佛里吉亚斗篷；她的敬意并不稍减，}
\verseline{485}{让他满载这些织物的礼物，又这样说道：}
\verseline{}{‘孩子，收下这些，让它们成为我双手留给你的纪念，}
\verseline{}{证明赫克托耳之妻安德罗玛刻}
\verseline{}{对你长久的爱。收下亲人最后的礼物，}
\verseline{}{你是我那阿斯提阿那克斯仅存的影像：}
\verseline{490}{他也曾有这样的眼睛、这样的手、这样的面容；}
\verseline{}{如今若还活着，应当正与你同岁，走向成年。’\footnote{阿斯提阿那克斯是赫克托耳与安德罗玛刻之子，特洛伊陷落后被杀。安德罗玛刻在阿斯卡尼乌斯的眼睛、手与面容中看见亡子的相似，并说他若长大，两人本应同龄。赠衣因此也寄托了对亡子的怀念。}}
\verseline{}{离开时，我眼泪涌出，对他们说：}
\verseline{}{‘幸福地生活吧，你们自己的命运}
\verseline{}{已经完成；我们却被一重命运接一重命运地召去。}
\verseline{495}{你们的安宁已经赢得；不必再耕犁海面，}
\verseline{}{不必再寻找那总向后退去的奥索尼亚}
\verseline{}{田野。你们眼前已有克珊托斯的复制和一座特洛伊}
\verseline{}{出自自己的双手。我祈愿它得到更好的}
\verseline{}{征兆，也不那么暴露在希腊人面前。}
\verseline{500}{倘若有一天我进入台伯河和邻近台伯河的田野，}
\verseline{}{看见赐给我族人的城墙，}
\verseline{}{我们终将让这两座血缘相连的城市和邻近的民族，}
\verseline{}{伊庇鲁斯与赫斯珀里亚——同一位达耳达诺斯是他们的始祖}
\verseline{}{同一种遭际属于他们——在心灵中合成}
\verseline{505}{同一个特洛伊。愿我们的后代也把此事放在心上。’}
\verseline{}{我们驶上大海，沿近旁的刻劳尼亚群山前行，}
\verseline{}{从这里渡往意大利，水路最短。}
\verseline{}{其间太阳沉落，幽暗的山峦罩上阴影；}
\verseline{}{我们在渴望已久的土地怀中、临着水边躺下，}
\verseline{510}{桨位已抽签分定，我们各自在干燥的海岸}
\verseline{}{照料身体；睡眠滋润疲惫的四肢。}
\verseline{}{夜在时辰的驱策下还未走到天穹中央：}
\verseline{}{警醒的帕利努鲁斯便从卧席起身，逐一}
\verseline{}{探察风向，侧耳捕捉空中的动静；}
\verseline{515}{注视无声天宇间滑行的一切星辰，}
\verseline{}{大角星、多雨的许阿得斯星群与大小二熊，}
\verseline{}{又转身观察佩着黄金武装的俄里翁。}
\verseline{}{看见晴朗天空下万象稳定，}
\verseline{}{他便从船尾发出清楚的信号；我们拔营，}
\verseline{520}{试探航路，展开船帆的羽翼。}
\verseline{}{群星退隐，曙光已经染红天空，}
\verseline{}{这时我们远远看见昏暗的山丘和低低的}
\verseline{}{意大利。‘意大利！’阿卡忒斯最先呼喊，}
\verseline{}{‘意大利！’同伴以欢乐的声音向它致意。\footnote{阿卡忒斯首先高呼“意大利”，同伴随即复诵；第523—524行三次出现 Italia。这是特洛伊人流亡后第一次实际看见意大利，随后四匹白马的征兆又预示战争。}}
\verseline{525}{父亲安喀塞斯给巨大的调酒钵戴上花冠}
\verseline{}{斟满纯酒，又呼唤众神，}
\verseline{}{站在最前方的船上：}
\verseline{}{‘掌握海洋、陆地与风暴的强大神明，}
\verseline{}{请以风为我们开出平坦的航道，吹送顺风。’}
\verseline{530}{期盼的风势渐强，港湾已经敞开}
\verseline{}{越来越近，弥涅耳瓦的神殿在堡顶显现。}
\verseline{}{同伴收起船帆，把船头转向海岸。}
\verseline{}{港口被东来的海浪弯成一道弓，}
\verseline{}{对面的礁石溅起咸涩的白沫；}
\verseline{535}{港湾本身隐藏在内；塔楼般的岩石垂下双臂，}
\verseline{}{成为两道壁垒，神殿却从海岸向后退去。}
\verseline{}{我在这里看见最初的征兆：四匹马，}
\verseline{}{正在辽阔的草原吃草，白得像雪。}
\verseline{}{父亲安喀塞斯说：‘好客的土地啊，你带来战争}
\verseline{540}{马为战争披挂，这些马群预示着战争的威胁。\footnote{安喀塞斯先以马为战争装备、披挂（bello armantur）而断为战争之兆；又说同一群马惯于套车，承受同轭的缰绳，因此也有和平希望。}}
\verseline{}{但这些四足兽也惯于套车}
\verseline{}{在同一副轭下和谐地忍受缰绳；}
\verseline{}{所以也有和平的希望。’于是我们敬告}
\verseline{}{披甲铿鸣的帕拉斯，她最先接纳了欢欣的我们，}
\verseline{545}{又在祭坛前用佛里吉亚长袍遮住头顶；}
\verseline{}{依照赫勒努斯当作首要之事交代的训示，}
\verseline{}{如仪向阿尔戈斯的朱诺焚献受命应献的祭礼。\footnote{特洛伊人在第545—547行按赫勒努斯的指令遮头献祭，并敬拜阿尔戈斯的朱诺；这实行了他先前关于祭仪和安抚朱诺的嘱咐。}}
\verseline{}{不再迟延，一切誓愿依次完成，}
\verseline{}{我们随即转过那些披着风帆的帆桁两端，}
\verseline{550}{离开希腊人的屋宇和可疑的土地。}
\verseline{}{从这里可以看见赫尔库勒斯的塔伦图姆海湾——倘若传说不假}
\verseline{}{对面，拉喀尼亚女神的神殿高高升起，\footnote{“拉喀尼亚女神”是拉喀尼翁海角的朱诺，此处借女神称号指她的神庙。第551行另写塔伦图姆海湾；“赫尔库勒斯的”修饰塔伦图姆，联系其建城传说。}}
\verseline{}{又有考洛尼亚的城堡和使船破碎的斯库拉刻乌姆。}
\verseline{}{随后，特里纳克里亚的埃特纳从远处海面显现，}
\verseline{555}{我们遥遥听见海的巨大呻吟和受浪击打的礁石}
\verseline{}{破碎的轰声传到岸边，}
\verseline{}{浅滩翻涌，沙粒与潮水混在一处。}
\verseline{}{父亲安喀塞斯说：‘这无疑就是那卡律布狄斯：}
\verseline{}{赫勒努斯曾预言这些暗礁，这些骇人的岩石。}
\verseline{560}{同伴们，逃开这里，一齐俯身奋桨！’}
\verseline{}{他们丝毫不怠地遵照命令，帕利努鲁斯首先}
\verseline{}{把轰鸣的船头转向左方的海浪。}
\verseline{}{全体船员也用桨与风帆驶向左侧。}
\verseline{}{弯曲的漩涡把我们举上天空；我们又}
\verseline{565}{在浪头撤去时，沉落到最深的亡灵之境。}
\verseline{}{礁石三次在空洞的岩隙间发出巨响：}
\verseline{}{我们三次看见白沫迸裂，水雾沾湿星辰。}
\verseline{}{其间，风和太阳把疲惫的我们一同抛下，}
\verseline{}{我们不识道路，滑向独眼巨人的海岸。}
\verseline{570}{这里的港湾广大，风无法侵入，}
\verseline{}{但埃特纳就在近旁，以骇人的崩裂轰鸣；}
\verseline{}{它时而向天空喷出一团黑云，}
\verseline{}{裹着沥青般的旋烟与灼亮的火山灰，}
\verseline{}{又举起团团火焰，舔舐星辰；}
\verseline{575}{时而吐出礁石与山体撕裂的内脏，}
\verseline{}{把它们抛向高空，将熔化的岩块伴着呻吟}
\verseline{}{卷作一团，深底也沸腾起来。}
\verseline{}{相传恩刻拉多斯被雷电烧焦半边的身躯}
\verseline{}{正被这座山体的重压镇住，巨大的埃特纳覆在他身上，}
\verseline{580}{从破裂的炉口呼出火焰；}
\verseline{}{每当他疲惫地翻转身躯，整个特里纳克里亚}
\verseline{}{便随轰鸣颤抖，烟尘在天空下织成帷幕。}
\verseline{}{那一夜，我们躲在树林中忍受这些庞然异象}
\verseline{}{却看不见是什么原因制造了声响。}
\verseline{585}{既没有星辰的火焰，也不见星光明亮的天穹；}
\verseline{}{只有乌云笼罩昏暗的天空，}
\verseline{}{无尽的深夜把月亮困在暴雨中。}
\verseline{}{第二日已经从东方最初的曙光升起，}
\verseline{}{奥罗拉从天穹移开湿润的暗影，}
\verseline{590}{突然，林中走出一个瘦到形销骨立的人，}
\verseline{}{陌生的形貌前所未见，衣着令人怜悯，}
\verseline{}{他向前走来，作求告的姿态，朝海岸伸出双手。}
\verseline{}{我们回望：可怕的污垢，蓬乱垂下的长须，}
\verseline{}{用荆刺缀合的衣服；但其余一切表明他是希腊人，}
\verseline{595}{也曾穿着祖国的铠甲，被派往特洛伊。}
\verseline{}{他远远看见达耳达诺斯人的服饰和特洛伊的}
\verseline{}{武装，被眼前景象吓得一时僵立，}
\verseline{}{停住脚步；不久又不顾一切地奔向海岸}
\verseline{}{带着泪水和祈求：‘我以星辰作证，}
\verseline{600}{以天上的众神和这可供呼吸的天空之光作证，}
\verseline{}{带走我吧，特乌克罗斯人，把我带到任何土地；}
\verseline{}{这就足够。我知道自己是达那俄斯舰队的一员，}
\verseline{}{也承认曾在战争中攻击伊利乌姆的家神；}
\verseline{}{如果我的罪所造成的伤害如此深重，}
\verseline{605}{就把我抛进波涛，沉入浩瀚的大海。}
\verseline{}{若我必死，死在人手中也使我甘愿。’}
\verseline{}{他说完便抱住我们的膝头，伏地辗转}
\verseline{}{紧紧不放。我们劝他说明自己是谁、出身何种血脉，}
\verseline{}{然后坦白什么命运正在驱赶他。}
\verseline{610}{父亲安喀塞斯本人没有多作迟疑，}
\verseline{}{便向青年伸出右手，以这一握安定他的心。}
\verseline{}{他终于放下恐惧，开口说道：}
\verseline{}{‘我的祖国是伊塔卡，我是不幸的乌利塞斯的同伴，}
\verseline{}{名叫阿开墨尼得斯；由于父亲阿达马斯托斯贫穷，我才}
\verseline{615}{前往特洛伊——但愿贫穷始终是我的命运！}
\verseline{}{同伴们惊惶地离开那残酷的门槛时，}
\verseline{}{竟忘了我，把我丢在独眼巨人}
\verseline{}{巨大的洞窟。那屋宇满是污血与血腥的筵席，}
\verseline{}{内部阴暗、庞大；他本人高耸，头颅撞击}
\verseline{620}{星辰——众神啊，从大地移开这样的灾害！——}
\verseline{}{没有人能轻易直视他的面孔，也不能同他交谈。}
\verseline{}{他吃的是可怜人的内脏与黑血。}
\verseline{}{我亲眼看见他以巨手抓起我们中的两人}
\verseline{}{他仰身坐在洞穴中央，将他们}
\verseline{625}{摔碎在岩石上；飞溅的污血漫过}
\verseline{}{门槛。我看见他咀嚼流淌黑色腐血的肢体}
\verseline{}{温热的关节在牙齿下颤抖。}
\verseline{}{但他并非不受惩罚；乌利塞斯没有忍下这种暴行，}
\verseline{}{伊塔卡人在巨大的危机中也不曾忘记自己。\footnote{阿开墨尼得斯是乌利塞斯遗落在独眼巨人岛上的希腊同伴。埃涅阿斯接纳昔日敌军的求告者；随后由他指认波吕斐摩斯，使《奥德赛》的旧故事从被遗弃者与特洛伊人的视角重新展开。}}
\verseline{630}{他吃饱食物，又被酒埋入昏睡，}
\verseline{}{歪着脖颈，巨大的身躯铺满洞穴，}
\verseline{}{睡梦中吐出污血与肉块，同血红的}
\verseline{}{酒混在一起；我们向伟大的神明祈祷}
\verseline{}{抽签分定各自的任务，从四周一齐}
\verseline{635}{涌上去，用削尖的木桩钻进他的眼睛——}
\verseline{}{那只巨眼藏在狰狞额头之下，独一无二，}
\verseline{}{大小如阿尔戈斯人的盾，或福玻斯的日轮——}
\verseline{}{终于欢喜地为同伴的亡魂复仇。}
\verseline{}{但可怜的人们，逃吧，快逃，把岸边的缆绳}
\verseline{640}{扯断！}
\verseline{}{和洞窟里圈养绵羊、挤取乳汁的}
\verseline{}{波吕斐摩斯一般形貌、一般庞大的}
\verseline{}{还有百名可怖的独眼巨人，散居在这片}
\verseline{}{弯曲的海岸，在高山间游荡。}
\verseline{645}{月亮的双角如今第三次被光辉填满，}
\verseline{}{我一直在林中，在野兽荒弃的}
\verseline{}{巢穴与洞窟之间拖延生命，从岩上眺望庞大的独眼巨人}
\verseline{}{一听见他们的脚步和声音就颤抖。}
\verseline{}{树枝给我悲惨的食物：浆果与坚硬的山茱萸果，}
\verseline{650}{拔起的草根喂养着我。}
\verseline{}{我扫视四方，这是我最先看见的舰队，}
\verseline{}{驶近这片海岸。无论它属于什么人，我已经}
\verseline{}{把自己交给它；只要逃过这可憎的种族就够了。}
\verseline{}{你们倒不如用任何一种死亡取走我的生命吧。’}
\verseline{655}{他刚刚说完，我们便看见山顶上}
\verseline{}{牧人波吕斐摩斯在羊群中移动庞大的}
\verseline{}{身躯，走向熟悉的海岸，}
\verseline{}{一个骇人、丑陋、巨大的怪物，眼里的光已经被夺去。}
\verseline{}{一根削去枝桠的松干引着他的手，也稳住脚步；}
\verseline{660}{披毛的绵羊一路伴随——那是他唯一的欢愉，}
\verseline{}{也是不幸中的慰藉。}
\verseline{}{他碰到深水，走进海面以后，}
\verseline{}{便洗去从被凿开的眼窝里流出的鲜血，}
\verseline{}{一面呻吟，一面咬响牙齿，涉过大海}
\verseline{665}{已经走到海水中央，浪涛仍未沾湿他高耸的腰腹。}
\verseline{}{我们惊惶地加速逃离那里，接纳了}
\verseline{}{这个确实应当救下的求告者，又默默割断缆绳；}
\verseline{}{我们俯身竞桨，翻搅海面。}
\verseline{}{他听见了，循着声音转过脚步；}
\verseline{670}{但右手已经无法碰到我们，}
\verseline{}{追赶的步伐也不能同伊奥尼亚的波浪匹敌，}
\verseline{}{他便发出无边的吼声，大海和所有}
\verseline{}{浪涛为之震颤，意大利的土地在深处}
\verseline{}{也惊恐起来，埃特纳弯曲的洞穴随之轰鸣。}
\verseline{675}{独眼巨人的族类被这吼声从树林和高山}
\verseline{}{惊起，纷纷奔向港湾，挤满海岸。}
\verseline{}{我们看见埃特纳的兄弟徒然站在那里，}
\verseline{}{以凶狠的眼光凝望，头颅高举入云，}
\verseline{}{那是一场骇人的集会：仿佛树顶高耸的}
\verseline{680}{参天橡树或结满球果的柏树，}
\verseline{}{一同站在朱庇特的高林，或狄安娜的圣林。}
\verseline{}{尖锐的恐惧逼得我们仓促抖开帆索，}
\verseline{}{不管驶向何方，便把船帆张向顺风。}
\verseline{}{可赫勒努斯的训示警告，不可把航线保持在斯库拉}
\verseline{685}{与卡律布狄斯之间；两边都离死亡只差一线，}
\verseline{}{于是我们决意让帆布向后，沿原路退去。}
\verseline{}{看啊，北风这时从佩洛鲁斯的狭窄关口}
\verseline{}{吹来。我驶过潘塔吉亚河天然岩石的出口，}
\verseline{}{墨伽拉的海湾和低伏的塔普索斯。}
\verseline{690}{不幸的乌利塞斯之同伴阿开墨尼得斯，}
\verseline{}{这样一处处指点——我们正倒着重走他曾迷失的海岸。\footnote{阿开墨尼得斯为新同伴指点海岸，舰队等于逆向重走乌利塞斯的迷航路线；一名希腊幸存者由此成为特洛伊人的向导，双方在独眼巨人的威胁下共同逃离。}}
\verseline{}{一座岛横在西坎尼亚海湾前方，}
\verseline{}{面对浪涛汹涌的普勒密里乌姆；古人称它}
\verseline{}{俄耳提癸亚。传说厄利斯的阿尔甫斯河\footnote{这里的俄耳提癸亚在西西里叙拉古附近；传统说阿尔甫斯河潜流海下，与阿瑞图萨泉相通。本卷75行也用俄耳提癸亚称提洛岛，两处同名而地点不同。}}
\verseline{695}{在海底开出隐秘的道路来到这里，如今}
\verseline{}{阿瑞图萨啊，它从你的泉口汇入西西里的波涛。}
\verseline{}{我们遵照命令，敬拜当地伟大的神明；随后，}
\verseline{}{我越过水漫成泽的赫洛鲁斯极其肥沃的土地。}
\verseline{}{接着擦过帕基努斯高耸的悬崖和前突的礁石}
\verseline{700}{远方出现命运从不容许人移动的}
\verseline{}{卡墨里那和革拉的平原，\footnote{卡墨里那是西西里城名，也是附近沼泽的名字。古代故事说神谕不许排干这片沼泽；居民排干后，敌军得以经此攻城。这里的“不容许人移动”承接这一禁令。}}
\verseline{}{以及取河流之名的巨城革拉。}
\verseline{}{然后高耸的阿克拉加斯从远处展示它宏大的}
\verseline{}{城墙，那里从前繁育雄骏的马匹；}
\verseline{705}{借着送来的风，我离开长满棕榈的塞利努斯，}
\verseline{}{又经过利吕贝乌姆礁石暗藏的危险浅滩。}
\verseline{}{德瑞帕努姆的港湾和令人悲痛的海岸，}
\verseline{}{就这样接纳了我。在这里，被大海无数风暴驱赶之后，}
\verseline{}{唉，我失去了父亲安喀塞斯——我在一切忧虑}
\verseline{710}{与厄运中的慰藉。在这里，最好的父亲啊，你把疲惫的我}
\verseline{}{抛下；唉，从那么多危险中被救出，竟全然徒劳！}
\verseline{}{预言者赫勒努斯向我警告许多骇人之事，}
\verseline{}{却未曾预言这场悲哀，可怕的刻莱诺也未曾预言。}
\verseline{}{这是最后一场苦难，也是漫长旅途的终点。}
\verseline{715}{我从那里离开，神明便把我送到你们的海岸。”\footnote{德瑞帕努姆是安喀塞斯死亡之地。第三卷以亡父之死结束，叙述随即回到第一卷狄多宫中；第五卷舰队重返西西里，并以周年祭赛公开纪念安喀塞斯。}}
\verseline{}{父亲埃涅阿斯就这样独自讲述，众人凝神倾听，}
\verseline{}{重述众神定下的命运，说明自己走过的航程。}
\verseline{}{终于，他沉默下来，在这里结束讲述，就此停歇。}
